\section{Ergebnisse}
\label{sec:results}

\subsection{Kern-Genauigkeiten}
\label{sec:accuracies}

Tabelle~\ref{tab:accuracy} fasst die FP32-, PTQ- und QAT-Genauigkeit für alle sechs Architektur$\times$Datensatz-Kombinationen zusammen (Quelle: \texttt{pipeline\_summary.csv}, \texttt{ptq\_summary.csv}, \texttt{qat\_summary.csv}). PTQ verursacht durchweg einen Genauigkeitsverlust gegenüber FP32, der auf ImageNet100 deutlich größer ausfällt als auf CIFAR10 (etwa \texttt{cnn}: $-17.98$ vs.\ $-10.22$ Prozentpunkte). QAT gleicht diesen Verlust nicht nur aus, sondern übertrifft in jeder Kombination sogar die FP32-Baseline, am deutlichsten beim \texttt{cnn} auf ImageNet100 (FP32 $43.04\,\%$, QAT $51.18\,\%$, $+8.14$ Punkte gegenüber FP32, $+26.12$ Punkte gegenüber PTQ). Dieser Effekt ist konsistent mit dem simulierten Quantisierungsrauschen als impliziter Regularisierung während des Trainings.

\begin{table*}[t]
\centering
\small
\caption{FP32-, PTQ- und QAT-Validierungsgenauigkeit (\%) für alle sechs Kombinationen, mit PTQ-Verlust und QAT-Verbesserung gegenüber PTQ (Prozentpunkte). Quelle: \texttt{pipeline\_summary.csv}.}
\label{tab:accuracy}
\begin{tabular}{llrrrrr}
\toprule
Datensatz & Modell & FP32 & PTQ & QAT & PTQ-Verlust & QAT-Erholung \\
\midrule
CIFAR10 & CNN & 81.30 & 71.08 & 84.46 & $-10.22$ & $+13.38$ \\
CIFAR10 & ResNet-18 & 84.18 & 82.76 & 90.51 & $-1.42$ & $+7.75$ \\
CIFAR10 & ResNet-50 & 81.14 & 79.54 & 86.80 & $-1.60$ & $+7.26$ \\
ImageNet100 & CNN & 43.04 & 25.06 & 51.18 & $-17.98$ & $+26.12$ \\
ImageNet100 & ResNet-18 & 80.06 & 71.30 & 82.52 & $-8.76$ & $+11.22$ \\
ImageNet100 & ResNet-50 & 80.02 & 64.06 & 83.22 & $-15.96$ & $+19.16$ \\
\bottomrule
\end{tabular}
\end{table*}

Auffällig ist, dass PTQ auf CIFAR10 (kleine $32\times32$-Eingaben) für die ResNets moderat ausfällt ($-1.4$ bis $-1.6$ Punkte), auf ImageNet100 aber massiv ($-8.8$ bis $-16.0$ Punkte). Diese Asymmetrie deutet bereits methodisch an, dass Datensatz und Auflösung die Quantisierungsempfindlichkeit stark mitbestimmen, ein Muster, das sich durch die gesamte Analyse zieht (\S\ref{sec:phase4}).

\subsection{Phase 1: Gewichts- vs.\ Aktivierungs-Dekomposition}
\label{sec:phase1-results}

Tabelle~\ref{tab:decomposition} zerlegt den PTQ- und QAT-Genauigkeitsverlust in einen Gewichts- und einen Aktivierungsanteil (Quelle: \texttt{activation\_decomposition.csv}; ausschließlich für CIFAR10 verfügbar). Unter PTQ dominiert in allen drei Architekturen der Gewichtsanteil den Genauigkeitsverlust deutlich, am stärksten beim \texttt{cnn} ($10.24$ gegenüber $-0.07$ Punkten), am moderatesten bei ResNet-50 ($1.90$ gegenüber $-0.23$ Punkten). Unter QAT kehrt sich das Bild um. Der isolierte Aktivierungspfad verliert relativ stärker, während der isolierte Gewichtspfad die FP32-Baseline übertrifft (negativer \enquote{Verlust}, das heißt eine Verbesserung), konsistent mit der QAT-Verbesserung aus Tabelle~\ref{tab:accuracy}.

\begin{table}[t]
\centering
\small
\caption{Gewichts- vs.\ Aktivierungs-Dekomposition, CIFAR10 (Prozentpunkte gegenüber FP32; negativ = Verbesserung). Quelle: \texttt{activation\_decomposition.csv}.}
\label{tab:decomposition}
\begin{tabular}{llrrl}
\toprule
Modell & Stufe & Gewicht & Akt.\ & dominant \\
\midrule
CNN & PTQ & 10.24 & $-0.07$ & Gewicht \\
CNN & QAT & $-3.22$ & 3.11 & Aktivierung \\
ResNet-18 & PTQ & 1.41 & 0.09 & Gewicht \\
ResNet-18 & QAT & $-6.36$ & $-5.48$ & Aktivierung \\
ResNet-50 & PTQ & 1.90 & $-0.23$ & Gewicht \\
ResNet-50 & QAT & $-5.50$ & $-3.72$ & Aktivierung \\
\bottomrule
\end{tabular}
\end{table}

\textbf{Konsequenz für die Hessian-Analyse.} Da die Hessian-basierte Sensitivitätsanalyse (\S\ref{sec:isolation}--\ref{sec:scores}) explizit nur Gewichtsschaden erfasst, ist ihr Anwendungsbereich durch diese Dekomposition abgesteckt. Unter PTQ, wo die Gewichte den Verlust dominieren, adressiert die Hessian-Analyse tatsächlich den Hauptmechanismus. Eine Übertragung auf die QAT-Sensitivität wäre dagegen fragwürdig, da dort die Aktivierungen relativ stärker zum Verlust beitragen. Dies ist ein Grund, warum die Hauptanalyse auf PTQ fokussiert wird (F3, teilweise beantwortet; ImageNet100-Daten zu dieser Dekomposition liegen nicht vor, siehe \S\ref{sec:limitations}).
\par\medskip

Ergänzend zeigt die Aktivierungs-Ablationsanalyse (\texttt{activation\_ablation.csv}), dass das sukzessive Zurücksetzen der Schichten mit dem höchsten Aktivierungs-Ausreißerfaktor (bis zu Faktor $25.1\times$ bei ResNet-50-\texttt{layer4.1.conv3}) auf FP32 nur geringe Genauigkeit zurückgewinnt (etwa $+0.16$ Punkte für ResNet-50/PTQ nach Ausschluss der fünf am stärksten ausreißenden Schichten). Aktivierungs-Ausreißer sind zwar messbar vorhanden (Wertebereichsfaktor \texttt{outlier\_factor} typischerweise $3$--$5\times$ bei Faltungsschichten, \texttt{activation\_ranges.csv}), tragen aber unter PTQ nicht dominant zum Gesamtschaden bei, konsistent mit dem in Tabelle~\ref{tab:decomposition} beobachteten \enquote{dominant=Gewicht}-Befund für PTQ.

\subsection{Phase 2: Random-Init-Kontrolle}
\label{sec:phase2}

Tabelle~\ref{tab:randominit} vergleicht die Trace der \texttt{conv1}-Schicht bei zufälliger Initialisierung (Mittelwert$\pm$Standardabweichung über drei Init-Seeds) mit der Trace desselben Layouts nach Training (FP32, unfusionierte Basis; Quelle: \texttt{random\_init\_comparison.csv}, \texttt{random\_init\_summary.csv}).

\begin{table}[t]
\centering
\small
\caption{\texttt{conv1}-Trace bei Zufallsinitialisierung vs.\ trainiert (FP32, unfusioniert). Klassifikation laut \texttt{random\_init\_comparison.csv}.}
\label{tab:randominit}
\begin{tabular}{lrrl}
\toprule
Modell & Random ($\mu\pm\sigma$) & Trainiert & Klasse \\
\midrule
CNN & $0.0024\pm0.0001$ & $65.87$ & gelernt \\
ResNet-18 & $7.22\pm1.19$ & $17.35$ & gelernt \\
ResNet-50 & $74.4\pm68.6$ & $15.55$ & gelernt \\
\bottomrule
\end{tabular}
\end{table}

In allen drei Modellen ist \texttt{conv1} als \enquote{gelernt} klassifiziert. Beim \texttt{cnn} liegt die trainierte Trace um den Faktor $\sim\!27\,500$ über der Zufalls-Trace, was die extremste der drei Verschiebungen darstellt. Bei ResNet-18 beträgt der Faktor $\sim\!2.4$, und bei ResNet-50 liegt die trainierte Trace unterhalb des mit sehr hoher Streuung behafteten Zufalls-Mittelwerts, wird aber dennoch als \enquote{gelernt} eingestuft, da die Elevation relativ zum medianen Trace-Wert des jeweiligen Modells nach Training deutlich zunimmt (\texttt{elev\_over\_median\_random}$=0.10$ vs.\ \texttt{elev\_over\_median\_trained}$=1.02$). Die Kernaussage ist über alle drei Architekturen konsistent. Die charakteristische \texttt{conv1}-Krümmungsspitze, die für die Prädiktor-Analyse in \S\ref{sec:phase4} zentral ist, entsteht durch Training und ist \emph{nicht} bereits eine Eigenschaft der Netzarchitektur bei Initialisierung. Für ResNet-50 zeigt der volle Trace-Profil-Vergleich (alle Schichten, nicht nur \texttt{conv1}) zusätzlich eine signifikante Spearman-Korrelation zwischen Random- und Trainiert-Profil ($\rho=0.60$, $p<0.001$) mit dem stärksten Ausreißer bei \texttt{layer4.1.conv2} (Elevation Random$=38.8\times$, Trainiert$=7.8\times$). Dies ist ein Hinweis darauf, dass bei ResNet-50 zusätzlich zur \texttt{conv1}-Dynamik auch architektonische, bereits bei Init vorhandene Effekte tiefer im Netz eine Rolle spielen, während beim ResNet-18-Profil keine signifikante Korrelation besteht ($\rho=0.09$, $p=0.70$).
\par\medskip

Diese tiefenabhängige Struktur bei Zufallsinitialisierung ist qualitativ konsistent mit der Mean-Field-Vorhersage von \citet{karakida_fisher_2019} (\S\ref{sec:fund-rmt}). Das Training kehrt dieses anfängliche Tiefenprofil für \texttt{conv1} systematisch um, indem es dessen Krümmung relativ zum Rest des Netzes anhebt.

\subsection{Phase 3: Trace-Drift und Rekonziliation}
\label{sec:phase3}

Im Projektverlauf entstanden für dieselbe nominelle Größe, die PTQ-Trace von ResNet-50-\texttt{conv1}, deutlich unterschiedliche historische Werte, abhängig davon, welche nicht dokumentierte Konfiguration ein früherer Lauf verwendete, nämlich $13\,161.8$ in einem Lauf und $1030.9$ in einem späteren. Die kanonische Neu-Messung (\S\ref{sec:canonical-config}) ergibt $97.76 \pm 6.27$ (Mittelwert$\pm$Standardabweichung über drei Proben-Seeds; \texttt{canonical\_traces.csv}), was Verhältnissen von $134.6\times$ bzw.\ $10.5\times$ zu den historischen Werten entspricht.

\begin{table*}[t]
\centering
\small
\caption{Rekonziliations-Ledger für ResNet-50-\texttt{conv1}: historische Werte gegen die kanonische Neu-Messung. Quelle: \texttt{trace\_reconciliation\_ledger.csv}.}
\label{tab:reconciliation}
\begin{tabularx}{\textwidth}{Xrrrl}
\toprule
Größe & Alter Wert & Kanonisch & Verhältnis & Erklärung \\
\midrule
\texttt{conv1} FP32 (Lauf 1) & 13.87 & 14.90 & $0.93\times$ & unfusionierte Basis (Residuum 2.8\,\%) \\
\texttt{conv1} FP32 (Lauf 2) & 11.79 & 14.90 & $0.79\times$ & ungeklärt \\
\texttt{conv1} PTQ (Lauf 1) & 13161.8 & 97.76 & $134.6\times$ & ungeklärt \\
\texttt{conv1} PTQ (Lauf 2) & 1030.9 & 97.76 & $10.5\times$ & ungeklärt \\
\texttt{conv1} Elevation FP32 & 14.0 & 11.73 & $1.19\times$ & ungeklärt \\
\bottomrule
\end{tabularx}
\end{table*}

Um diese Diskrepanzen zu erklären, wird ein systematischer Sweep über Basis (fusioniert/unfusioniert), Loss-Reduktion (Mittelwert/Summe), Proben-Anzahl, Datenquelle und Recheneinheit durchgeführt (\texttt{drift\_diagnosis.csv}). Nur der FP32-Wert aus Lauf~1 lässt sich reproduzieren, und zwar durch die unfusionierte Basis mit einem Residuum von $2.8\,\%$. Keine getestete Konfiguration reproduziert die beiden historischen PTQ-Werte auch nur annähernd. Selbst die extremsten getesteten Varianten, also die Summenreduktion ($14\,101$ für FP32, $7742$ für PTQ), das kleinste Proben-Budget (\texttt{maxIter}$=20$: $139.9$ bzw.\ $82.5$) oder ungeseedete Trainingsdaten mit Augmentation ($286.0$/$184.4$ bzw.\ $105.6$/$388.6$), bleiben um mindestens eine Größenordnung von $13\,161.8$ entfernt.

Dies wird als dokumentierter, abgeschlossener Befund berichtet und nicht als offene Frage. Die historischen PTQ-Werte lassen sich mit den verfügbaren Konfigurationsknöpfen nicht rekonstruieren, und plausible Ursachen sind ein damals nicht protokollierter Checkpoint-, Kalibrierungs- oder Codepfad-Unterschied, der außerhalb des hier gesweepten Parameterraums liegt. Die praktische Konsequenz ist die in \S\ref{sec:canonical-config} eingeführte Disziplin. Ohne eine explizit dokumentierte, fixierte Konfiguration sind Krümmungs-basierte Sensitivitätswerte zwischen Läufen, und potenziell zwischen Publikationen, nicht direkt vergleichbar.

\subsection{Phase 3.5: Fusion vs.\ Quantisierung}
\label{sec:phase35}

Batch-Normalisierungs-Fusion allein erhöht die \texttt{conv1}-Trace bereits erheblich, bevor überhaupt quantisiert wird (Verhältnis fusioniert/unfusioniert, FP32, \texttt{canonical\_traces.csv}): $11.7\times$ bei ResNet-50 ($174.5$ vs.\ $14.9$), $10.7\times$ bei ResNet-18 ($182.1$ vs.\ $17.0$), $7.3\times$ beim \texttt{cnn} ($460.5$ vs.\ $63.1$). Tabelle~\ref{tab:fusion} setzt dies in Bezug zur zusätzlichen Amplifikation durch PTQ selbst (\texttt{conv1\_ptq\_amplification} $=\mathrm{Tr}_{\mathrm{PTQ}}/\mathrm{Tr}_{\mathrm{FP32,fused}}$, sowie deren Rang unter allen Schichten des Modells; Quelle: \texttt{quant\_induced\_summary.csv}).

\begin{table*}[t]
\centering
\small
\caption{Fusion- und PTQ-Amplifikation für \texttt{conv1}. Rang = Position von \texttt{conv1} in der Rangfolge aller Schichten nach PTQ-Amplifikation ($1=$ am stärksten amplifiziert). Quelle: \texttt{quant\_induced\_summary.csv}.}
\label{tab:fusion}
\begin{tabular}{lrrrl}
\toprule
Modell & Fusion & PTQ-Amp. & Rang & Verdikt \\
\midrule
CNN & $7.3\times$ & $1.15$ & 6/6 & quant\_induced \\
ResNet-18 & $10.7\times$ & $0.95$ & 19/21 & fp32\_intrinsic \\
ResNet-50 & $11.7\times$ & $0.56$ & 51/54 & fp32\_intrinsic \\
\bottomrule
\end{tabular}
\end{table*}

Nur beim \texttt{cnn} übersteigt die PTQ-Trace die bereits fusionierte FP32-Trace nennenswert (Faktor $1.15$); bei ResNet-18 und ResNet-50 liegt die PTQ-Trace \emph{unterhalb} der fusionierten FP32-Trace (Faktoren $0.95$ und $0.56$), und in beiden Fällen liegt \texttt{conv1} in der PTQ-Amplifikations-Rangfolge nahe am unteren Ende (Rang 19 von 21 bzw.\ 51 von 54). Die dramatische \texttt{conv1}-Erhöhung, die in \S\ref{sec:phase2} beobachtet wurde, ist demnach primär ein Fusions-Effekt (bzw.\ ein Trainingseffekt, siehe Phase~2) und nicht primär ein Quantisierungs-induzierter Effekt, mit der Ausnahme des \texttt{cnn}, wo PTQ selbst eine zusätzliche, wenn auch moderate Verstärkung beiträgt. Diese Unterscheidung ist wichtig für die korrekte Interpretation von Tabelle~\ref{tab:rank} in \S\ref{sec:phase4}. Die hohe \texttt{conv1}-Krümmung, die dort als Prädiktor-Signal dient, entsteht überwiegend \emph{vor} der eigentlichen PoT-Rundung.

\subsection{Phase 4: Prädiktor-Test (Kernergebnis)}
\label{sec:phase4}

Dies ist die zentrale Analyse des Projekts und die Grundlage des Workshop-Papers. Tabelle~\ref{tab:rank} zeigt für alle sechs Kombinationen (PTQ, loss-basierter Schaden), welchen Rang $S_{\mathrm{raw}}$ der Schicht mit dem tatsächlich größten Schaden zuweist, sowie die top-$k^*$-Überlappung aller drei Scores (Quelle: \texttt{weight\_ablation\_loss\_damage\_correlation.csv}, \texttt{normalized\_ranks\_loss.csv}).

\begin{table*}[t]
\centering
\small
\caption{Rang des wahren Top-Damage-Layers unter $S_{\mathrm{raw}}$, dessen Spearman-$\rho$ mit 95\%-Bootstrap-CI, und top-$k^*$-Überlappung aller drei Scores, PTQ, loss-basiert. $k^*=1,3,6$ für CNN/ResNet-18/ResNet-50. Fett = CI schließt Null aus. Quelle: \texttt{normalized\_ranks\_loss.csv}, \texttt{bootstrap\_ci\_spearman\_loss.csv}.}
\label{tab:rank}
\begin{tabular}{llcccccc}
\toprule
Datensatz & Modell ($n$) & Top-Damage-Layer & $S_{\mathrm{raw}}$-Rang & $S_{\mathrm{raw}}\ \rho$ [95\%-CI] & $S_{\mathrm{raw}}$ top-$k^*$ & $S_{\mathrm{pert}}$ top-$k^*$ & $S_{\mathrm{hawq}}$ top-$k^*$ \\
\midrule
ImageNet100 & CNN (6) & \texttt{conv1} & 3 & $0.54\,[-1.00,0.80]$ & 0/1 & 0/1 & 0/1 \\
ImageNet100 & ResNet-18 (21) & \texttt{conv1} & \textbf{1} & $0.10\,[-0.43,0.54]$ & 1/3 & 0/3 & 0/3 \\
ImageNet100 & ResNet-50 (54) & \texttt{conv1} & \textbf{1} & $0.13\,[-0.17,0.41]$ & 1/6 & 0/6 & 0/6 \\
CIFAR10 & CNN (6) & \texttt{conv1} & 2 & $0.77\,[0.00,1.00]$ & 0/1 & 0/1 & 0/1 \\
CIFAR10 & ResNet-18 (21) & \texttt{layer1.1.conv1} & \textbf{1} & $0.42\,[-0.01,0.75]$ & 1/3 & 0/3 & 0/3 \\
CIFAR10 & ResNet-50 (54) & \texttt{fc} & 40 & $\mathbf{0.36\,[0.10,0.58]}$ & 2/6 & 1/6 & 1/6 \\
\bottomrule
\end{tabular}
\end{table*}

$S_{\mathrm{raw}}$ erreicht eine Rang-1-Identifikation in 3 von 6 Kombinationen und wird in der top-$k^*$-Überlappung von keinem der beiden anderen Scores übertroffen. In beiden \texttt{cnn}-Zeilen liegt es gleichauf (dort erlaubt $k^*=1$ keinem Score, die einzige wahre Top-Schicht zu treffen), in den übrigen vier Kombinationen ist es strikt überlegen. $S_{\mathrm{pert}}$ erreicht in keiner Kombination Rang~1 und versenkt die wahre Top-Schicht regelmäßig weit hinten in der Rangfolge, insbesondere wenn diese eine Faltungsschicht mit kleinem Fan-in ist (\texttt{conv1}, \texttt{layer1.1.conv1}). $S_{\mathrm{hawq}}$ erreicht ebenfalls nie Rang~1 und übertrifft $S_{\mathrm{raw}}$s top-$k^*$-Überlappung in keiner Kombination. ResNet-50/CIFAR10 ist der Fall, in dem $S_{\mathrm{raw}}$ selbst am stärksten versagt (Rang 40 von 54), konsistent damit, dass der Schaden dort am diffusesten über alle 54 Schichten verteilt ist. Hier erreicht keiner der drei Scores mehr als ein Drittel der wahren top-$k^*$-Menge. Von den sechs $S_{\mathrm{raw}}$-Korrelationen in der Tabelle schließt nur eine, ResNet-50/CIFAR10, das 95\%-Bootstrap-CI von Null aus; die übrigen fünf sind richtungskonsistent positiv, aber bei dieser Stichprobengröße ($n=6$ bis $54$ Schichten) statistisch nicht gesichert, ein Punkt, der in \S\ref{sec:phase4} unten weiter vertieft wird.

Abbildung~\ref{fig:rankcurve} bündelt diese sechs Einzelfälle zu einer über die Architekturen normalisierten Zusammenfassung. Für eine Rang-Schwelle $t\in[0,1]$ zeigt sie den Anteil der sechs PTQ-Kombinationen, in denen der jeweilige Score die wahre Top-Damage-Schicht innerhalb der obersten $t$-Schichten platziert. $S_{\mathrm{raw}}$ dominiert über den gesamten Schwellwertbereich, erreicht bereits bei $t=0.05$ eine Trefferquote von $3/6$ und bei $t=0.5$ von $5/6$, während $S_{\mathrm{pert}}$ selbst bei $t=0.5$ noch keine einzige Kombination erfasst (seine engste Platzierung liegt bei $t=0.31$) und $S_{\mathrm{hawq}}$ dazwischen liegt.

\begin{figure}[t]
\centering
\includegraphics[width=0.85\linewidth]{figures/fig_06_ranking_curves.pdf}
\caption{Kumulative Rang-Kurven: Anteil der 6 PTQ-Kombinationen, in denen die wahre Top-Damage-Schicht innerhalb der obersten $t$-Schichten (x-Achse) liegt, pro Score. Graue Diagonale = Zufalls-Rang-Erwartung. Erzeugt von \texttt{notebooks/16\_full\_ranking\_curves.ipynb}.}
\label{fig:rankcurve}
\end{figure}
\par\medskip

Abbildung~\ref{fig:main} zeigt die vollständige Score-Wert-vs-Schaden-Beziehung für alle sechs Kombinationen mit den zugehörigen Spearman-Korrelationen und Bootstrap-Konfidenzintervallen. Die volle Rangkorrelation ist insgesamt fragil. Von 18 Score$\times$Kombination-Paaren sind 5 bei $p<0.05$ signifikant (4, wenn ein Konfidenzintervall, das die Null ausschließt, als Kriterium verwendet wird), davon vier für $S_{\mathrm{pert}}$ bzw.\ $S_{\mathrm{hawq}}$ mit \emph{negativem} Vorzeichen. $S_{\mathrm{pert}}$ korreliert in 5 von 6 Kombinationen negativ mit dem Schaden, entgegengesetzt zu $S_{\mathrm{raw}}$s durchweg positivem Vorzeichen in denselben 5 Kombinationen.

\begin{figure}[t]
\centering
\includegraphics[width=\linewidth]{figures/fig_01_facet_plot.pdf}
\caption{Score-Wert (log-Skala) vs.\ gemessener \texttt{abs\_loss\_damage}, PTQ; obere Zeile ImageNet100, untere Zeile CIFAR10, Spalten CNN/ResNet-18/ResNet-50. Offene Diamanten markieren die wahre Top-Damage-Schicht (beschriftet, wo nicht \texttt{conv1}). Jedes Panel-Inset zeigt Spearman-$\rho$ mit 95\%-Bootstrap-CI ($B=2000$); fett = CI schließt Null aus. Erzeugt von \texttt{notebooks/06\_main\_facet\_plot.ipynb}.}
\label{fig:main}
\end{figure}

\textbf{Parameter-Normalisierung.} Eine parameter-normalisierte Variante von $S_{\mathrm{raw}}$ (Division durch die Parameterzahl der Schicht) verbessert die Korrelation in 5 von 6 Kombinationen (etwa \texttt{cnn}/CIFAR10: $\rho=0.77\rightarrow1.00$), verschlechtert sie aber bei ResNet-50/CIFAR10 deutlich ($\rho=0.36\rightarrow0.16$; Quelle: \texttt{raw\_trace\_per\_param\_comparison\_loss.csv}). Da kein Normalisierungsschema universell überlegen ist, bleibt die unnormalisierte Form das Hauptresultat (\S\ref{sec:scores}).
\par\medskip

\textbf{Metrik-Konkordanz.} Tabelle~\ref{tab:concordance} zeigt die paarweise Spearman-Korrelation zwischen den drei Score-Werten selbst (unabhängig vom Schaden-Ziel; Quelle: \texttt{predictor\_concordance.csv}). Bei ResNet-50 sind $S_{\mathrm{pert}}$ und $S_{\mathrm{hawq}}$ nahezu identisch ($\rho=0.84$ CIFAR10, $\rho=0.83$ ImageNet100), sodass der Drei-Wege-Vergleich dort effektiv zu einem Zwei-Wege-Vergleich wird ($S_{\mathrm{raw}}$ gegen ein fast identisches $S_{\mathrm{pert}}$/$S_{\mathrm{hawq}}$-Paar). Dass $S_{\mathrm{hawq}}$s top-$k^*$-Überlappung bei ResNet-50 (Tabelle~\ref{tab:rank}) $S_{\mathrm{raw}}$ nie übertrifft, ist somit größtenteils $S_{\mathrm{pert}}$s Muster, das durch das Produkt hindurchgereicht wird, und kein unabhängiges drittes Ergebnis.

\begin{table*}[t]
\centering
\small
\caption{Paarweise Spearman-$\rho$ zwischen den drei Score-Werten, PTQ. Quelle: \texttt{predictor\_concordance.csv}.}
\label{tab:concordance}
\begin{tabular}{llrrr}
\toprule
Modell & Datensatz & $\rho(S_{\mathrm{raw}},S_{\mathrm{pert}})$ & $\rho(S_{\mathrm{raw}},S_{\mathrm{hawq}})$ & $\rho(S_{\mathrm{pert}},S_{\mathrm{hawq}})$ \\
\midrule
CNN & ImageNet100 & $-0.71$ & $0.60$ & $0.09$ \\
CNN & CIFAR10 & $-0.54$ & $0.71$ & $0.03$ \\
ResNet-18 & ImageNet100 & $-0.60$ & $0.43$ & $0.41$ \\
ResNet-18 & CIFAR10 & $-0.55$ & $0.45$ & $0.43$ \\
ResNet-50 & ImageNet100 & $0.40$ & $0.78$ & $\mathbf{0.83}$ \\
ResNet-50 & CIFAR10 & $0.48$ & $0.86$ & $\mathbf{0.84}$ \\
\bottomrule
\end{tabular}
\end{table*}

Die vorstehende Analyse ist auf PTQ skaliert; der direkte Vergleich mit der QAT-Krümmungsstruktur, der F1 beantwortet, folgt in \S\ref{sec:phase9}.

\subsection{Phase 5: Spike-Selection}
\label{sec:phase5}

Tabelle~\ref{tab:spike} vergleicht für CIFAR10 die Schicht mit der höchsten Krümmung im trainierten FP32-Modell (\enquote{Trace-Spike}) mit der Schicht mit dem größten gemessenen Loss-Schaden (\enquote{Damage-Spike}), Quelle \texttt{spike\_selection.csv} (kanonische fusionierte Basis für die Trace-Spike-Bestimmung, konsistent mit \texttt{canonical\_traces.csv}).

\begin{table*}[t]
\centering
\small
\caption{Trace-Spike (höchste FP32-Trace, kanonische fusionierte Basis) vs.\ Damage-Spike (größter PTQ-Loss-Schaden), CIFAR10.}
\label{tab:spike}
\begin{tabular}{lllc}
\toprule
Modell & Trace-Spike & Damage-Spike & Übereinst.\ \\
\midrule
CNN & \texttt{conv2} & \texttt{conv1} & nein \\
ResNet-18 & \texttt{layer1.1.conv1} & \texttt{layer1.1.conv1} & \textbf{ja} \\
ResNet-50 & \texttt{conv1} & \texttt{fc} & nein \\
\bottomrule
\end{tabular}
\end{table*}

Unter der kanonischen fusionierten Basis stimmen Trace-Spike und Damage-Spike in 1 von 3 CIFAR10-Modellen überein, nämlich bei ResNet-18. Beim \texttt{cnn} liegt die Trace-Spitze bei \texttt{conv2}, während \texttt{conv1} am stärksten geschädigt wird; bei ResNet-50 liegt die Trace-Spitze bei \texttt{conv1} (Fusions-getrieben, \S\ref{sec:phase35}), während \texttt{fc} am stärksten geschädigt wird. Auf ImageNet100 (ebenfalls \texttt{spike\_selection.csv}) fällt die Übereinstimmung höher aus. ResNet-50 und ResNet-18 stimmen beide überein (\texttt{conv1} ist dort sowohl Trace- als auch Damage-Spike), nur das \texttt{cnn} weicht ab (Trace-Spitze \texttt{conv3}, Damage-Spitze \texttt{conv1}), womit 2 von 3 ImageNet100-Modellen übereinstimmen. Dieser Datensatz-Unterschied ist konsistent mit dem in \S\ref{sec:phase4} beobachteten Muster, dass \texttt{conv1} auf ImageNet100 durchweg die am stärksten geschädigte Schicht ist, auf CIFAR10 dagegen nicht immer.
\par\medskip

Bemerkenswert ist zusätzlich, dass diese Übereinstimmung selbst basisabhängig ist. Verwendet man statt der kanonischen fusionierten FP32-Trace die Trace-Elevation relativ zur Zufallsinitialisierung aus der Random-Init-Vergleichsanalyse (\texttt{random\_init\_summary.csv}, unfusionierte Basis, \S\ref{sec:phase2}), so ergibt sich für CIFAR10 eine andere Trace-Spike-Schicht pro Modell (ResNet-50: \texttt{layer4.1.conv2}; ResNet-18: \texttt{layer3.0.conv2}; \texttt{cnn}: \texttt{conv3}), und keines der drei Modelle stimmt mehr mit seinem Damage-Spike überein. Dieser Kontrast ist selbst ein Ergebnis. Ob die naive Heuristik \enquote{höchste Krümmung = größter Schaden} zutrifft, hängt nicht nur vom Modell, sondern auch von der gewählten Trace-Konfiguration ab, was ein weiterer Beleg für die in \S\ref{sec:phase3} dokumentierte Konfigurations-Sensitivität ist.

\subsection{Phase 6: Architektonische Ursachensuche}
\label{sec:phase6}

Für die jeweilige Trace-Spike-Schicht jedes Modells (\S\ref{sec:phase5}) wird versucht, die gemessene Pro-Parameter-Trace über die KFAC-Faktorisierung $\mathrm{Tr}(F)\approx\mathrm{Tr}(A)\cdot\mathrm{Tr}(G)$ vorherzusagen (\S\ref{sec:fund-rmt}). Tabelle~\ref{tab:kfac} zeigt, dass die Zerlegung in allen getesteten Fällen numerisch scheitert, da die vorhergesagte Pro-Parameter-Trace von der tatsächlich gemessenen um 3 bis über 8 Größenordnungen abweicht (Quelle: \texttt{spike\_layer\_kfac.csv}).

\begin{table*}[t]
\centering
\small
\caption{KFAC-Vorhersage vs.\ gemessene Pro-Parameter-Trace für die jeweilige Trace-Spike-Schicht, PTQ. Verhältnis $=$ gemessen / vorhergesagt. Quelle: \texttt{spike\_layer\_kfac.csv}.}
\label{tab:kfac}
\begin{tabular}{llrrr}
\toprule
Modell & Schicht & Vorhergesagt & Gemessen & Verhältnis \\
\midrule
CNN & \texttt{conv1} & $3.6\times10^{-9}$ & $0.614$ & $1.7\times10^{8}$ \\
CNN & \texttt{conv2} & $2.3\times10^{-9}$ & $0.069$ & $3.0\times10^{7}$ \\
ResNet-18 & \texttt{layer1.1.conv1} & $3.5\times10^{-11}$ & $0.0069$ & $2.0\times10^{8}$ \\
ResNet-50 & \texttt{conv1} & $3.7\times10^{-9}$ & $0.0566$ & $1.5\times10^{7}$ \\
ResNet-50 & \texttt{fc} & $1.3\times10^{-7}$ & $3.3\times10^{-4}$ & $2.5\times10^{3}$ \\
\bottomrule
\end{tabular}
\end{table*}

Die Abweichung ist für Faltungsschichten (\texttt{conv1}, \texttt{conv2}, \texttt{layer1.1.conv1}) systematisch größer (7--8 Größenordnungen) als für die vollverbundene \texttt{fc}-Schicht von ResNet-50 (gut 3 Größenordnungen). Dies ist ein Hinweis, dass die KFAC-Näherung, die ursprünglich für einzelne dichte Schichten entwickelt \citep{martens_kfac_2015} und erst später auf Faltungsschichten erweitert wurde \citep{grosse_kfc_2016}, für die hier betrachteten Faltungsschichten und deren spezifische Fan-in/Fan-out-Struktur (Tabelle~\ref{tab:descriptors}) besonders schlecht kalibriert ist. Auch die aus der Faktorisierung abgeleiteten $A$- und $G$-Anteile (\texttt{spike\_A\_share}, \texttt{spike\_G\_share}) sind inkonsistent. Sie liegen wiederholt außerhalb des sinnvollen Bereichs $[0,1]$ (etwa \texttt{cnn}/\texttt{conv1}: $A$-Anteil $16.6$, $G$-Anteil $-15.6$), was auf ein numerisch instabiles Verhältnis zwischen den Faktoren $\mathrm{Tr}(A)$ und $\mathrm{Tr}(G)$ hindeutet und nicht auf eine sinnvolle additive Zerlegung.

\begin{table}[t]
\centering
\footnotesize
\caption{Strukturelle Deskriptoren der Trace-Spike-Schichten. Quelle: \texttt{spike\_layer\_descriptors.csv}.}
\label{tab:descriptors}
\begin{tabular}{p{1.55cm}p{1.75cm}rrr}
\toprule
Modell & Schicht & Fan-in & Fan-out & \#Param.\ \\
\midrule
CNN & \texttt{conv1} & 27 & 288 & 864 \\
CNN & \texttt{conv2} & 288 & 576 & 18432 \\
ResNet-18 & \texttt{layer1.1.\allowbreak conv1} & 576 & 576 & 36864 \\
ResNet-50 & \texttt{conv1} & 27 & 576 & 1728 \\
ResNet-50 & \texttt{fc} & 2048 & 10 & 20480 \\
\bottomrule
\end{tabular}
\end{table}

Das KFAC-Scheitern wird als ehrliches Negativresultat gewertet. Die in \S\ref{sec:fund-rmt} skizzierte Hoffnung, Krümmungsspitzen über eine Eingabe-/Gradienten-Zerlegung architektonisch zu erklären, lässt sich mit der verwendeten Implementierung im PoT-Kontext nicht einlösen. Die drei Hypothesen H1--H3 zur architektonischen Ursache (Fan-in-getrieben, Output-Map-Größe-getrieben, vorangehender-$A$-Faktor-getrieben; \texttt{spike\_layer\_attribution.csv}) konnten entsprechend nicht diskriminiert werden. Die Spalte \texttt{hypothesis\_supported} ist für alle getesteten Schicht-Deskriptor-Paare \enquote{undetermined}, unter anderem weil der für eine Diskriminierung vorgesehene Auflösungskontrast zwischen CIFAR10 und ImageNet100 (\texttt{resolution\_contrast\_ratio}) für keine Kombination auswertbar war (durchweg \texttt{NaN}).

\subsection{QAT-Ergebnisse und PTQ-vs-QAT-Vergleich}
\label{sec:phase9}

Dieser Abschnitt beantwortet F1: Wie unterscheiden sich PTQ und QAT in der schichtweisen Verlust-Krümmung, und wie wirkt sich das auf die Score-vs-Schaden-Beziehung aus? Die AXIOM-Einreichung beschränkt sich auf PTQ; dieser Bericht ergänzt QAT auf beiden Ebenen.

\textbf{Die Krümmungsstruktur selbst ändert sich unter QAT kaum.} Vergleicht man die kanonische, fusionierte Hessian-Spur \emph{pro Schicht} zwischen dem PTQ- und dem QAT-Checkpoint derselben Architektur$\times$Datensatz-Kombination (Quelle: \texttt{weight\_ablation\_canonical\_v2.csv} und dessen ImageNet100-Variante), ergibt sich eine durchweg sehr hohe Spearman-Korrelation zwischen beiden Trace-Profilen ($\rho=0.95$--$1.00$ über alle sechs Kombinationen, $p<0.001$ in jedem Fall). Welche Schichten relativ zueinander stark oder schwach gekrümmt sind, ist also weitgehend unabhängig davon, ob das Netz per PTQ nachträglich oder per QAT von Beginn an quantisierungsbewusst trainiert wurde. Das ist konsistent mit dem Befund aus \S\ref{sec:phase35}, dass die dominanten Krümmungsspitzen primär architektonisch (Fusion, Fan-in) und nicht durch den spezifischen Trainingspfad bedingt sind.

\textbf{Die Score-vs-Schaden-Beziehung bricht unter QAT dagegen fast vollständig zusammen.} Abbildung~\ref{fig:qat} zeigt das QAT-Gegenstück zu Abbildung~\ref{fig:main} (Bootstrap-Rho und CI hier direkt aus den QAT-Per-Layer-Werten neu berechnet, da \texttt{bootstrap\_ci\_spearman\_loss.csv} nur PTQ abdeckt; Methodik identisch zu \S\ref{sec:scores}, $B=2000$, Seed 42). Von den 18 loss-basierten QAT-Score$\times$Kombination-Paaren ist \emph{keines} bei $p<0.05$ signifikant (kleinster $p$-Wert $0.10$, ResNet-18/CIFAR10, $S_{\mathrm{pert}}$); das ist sogar schwächer als das eine genauigkeitsbasierte QAT-Ergebnis, das im ursprünglichen Berichtsentwurf notiert wurde (ResNet-18/ImageNet100, $S_{\mathrm{pert}}$, $\rho=-0.47$, $p=0.031$; Quelle: \texttt{weight\_ablation\_canonical\_correlation\_v2.csv}).

\begin{figure}[t]
\centering
\includegraphics[width=\linewidth]{figures/fig_03_qat_facet.pdf}
\caption{QAT-Gegenstück zu Abbildung~\ref{fig:main}: Score-Wert (log-Skala) vs.\ \texttt{abs\_loss\_damage}, QAT. Bootstrap-CI direkt aus den QAT-Werten berechnet ($B=2000$, Seed 42). Erzeugt von \texttt{notebooks/14\_qat\_facet\_plot.ipynb}.}
\label{fig:qat}
\end{figure}

Der Grund liegt in \S\ref{sec:phase1-results}: Unter QAT verbessert die isolierte Quantisierung einer Schicht den Loss häufig gegenüber der FP32-Baseline (negativer Schaden), was Vorzeichen und Größenordnung des Zielwerts gegenüber PTQ grundlegend verändert. Die auf PTQ kalibrierte Score-Intuition, dass hohe Krümmung mit hohem Schaden einhergeht, überträgt sich damit nicht auf ein Regime, in dem \enquote{Schaden} teils gar keiner ist. Antwort auf F1: PTQ und QAT reformen die \emph{relative} Krümmungsstruktur eines Netzes kaum, sie reformen aber grundlegend, was diese Krümmung über den tatsächlichen Quantisierungseffekt aussagt.

\subsection{Signed vs.\ Absolute Damage}
\label{sec:signed}

Die Kernanalyse (\S\ref{sec:phase4}) rankt Schichten nach $\mathrm{abs\_loss\_damage}$, da ein reines Vorzeichen-Ranking Schäden beider Richtungen ungleich behandeln würde (\S\ref{sec:isolation}). Abbildung~\ref{fig:signed} zeigt zur Einordnung dasselbe PTQ-Raster wie Abbildung~\ref{fig:main}, jedoch mit dem \emph{vorzeichenbehafteten} $\mathrm{loss\_damage}$ auf der y-Achse.

\begin{figure}[t]
\centering
\includegraphics[width=\linewidth]{figures/fig_04_signed_damage.pdf}
\caption{Wie Abbildung~\ref{fig:main}, jedoch mit vorzeichenbehaftetem $\mathrm{loss\_damage}$ (PTQ) auf der y-Achse; gestrichelte Linie bei Null. Erzeugt von \texttt{notebooks/15\_signed\_damage\_plot.ipynb}.}
\label{fig:signed}
\end{figure}

Unter PTQ ist der Vorzeichenwechsel selten, aber nicht vernachlässigbar: von den 162 Schicht-Messungen über alle sechs Kombinationen sind 49 ($30\,\%$) negativ, das heißt, die isolierte PTQ-Quantisierung dieser Schichten senkt den Loss geringfügig unter die FP32-Baseline (Quelle: \texttt{weight\_ablation\_loss\_damage.csv}). Da diese negativen Werte betragsmäßig überwiegend klein sind, ändert sich die Identität der jeweils am stärksten geschädigten Schicht zwischen vorzeichenbehafteter und Betrags-Rangfolge in keiner der sechs Kombinationen. Der $|\cdot|$-Betrag ist damit für das Ranking die richtige Wahl (\S\ref{sec:isolation}), während die vorzeichenbehaftete Ansicht zusätzlich zeigt, dass \enquote{Schaden} unter PTQ keine Einbahnstraße ist, ein Muster, das unter QAT (\S\ref{sec:phase9}) noch deutlich ausgeprägter auftritt.

\subsection{Deployment-Ergebnisse}
\label{sec:phase10}

Tabelle~\ref{tab:deployment} zeigt die gemessene mittlere Laufzeit-Beschleunigung (INT8 relativ zu FP32, über eine Reihe von Batch-Größen gemittelt) auf einer NVIDIA A100-GPU für zwei torchao-0.17-Konfigurationen: die reine Gewichtsquantisierung (\texttt{int8\_weight\_only}) und die dynamische Aktivierungsquantisierung (\texttt{int8\_dynamic\_act}). Quelle ist ein projektinternes Diagnose-Protokoll (\texttt{int8\_perf\_diagnosis.txt}), das für jede der sechs Kombinationen und beide Quantisierungsstufen (PTQ/QAT-Checkpoint als Ausgangspunkt der Deployment-Konvertierung) einen vollständigen Batch-Größen-Sweep protokolliert.

\begin{table*}[t]
\centering
\small
\caption{Mittlere INT8-vs-FP32-Laufzeit-Beschleunigung auf NVIDIA A100 (torchao 0.17), gemittelt über Batch-Größen $1$--$32$. Werte $>1$ = schneller als FP32. Quelle: \texttt{int8\_perf\_diagnosis.txt}.}
\label{tab:deployment}
\begin{tabular}{llrr}
\toprule
Datensatz & Modell & wt.-only & dyn.-akt.\ \\
\midrule
CIFAR10 & CNN & $1.05$ & $0.21$ \\
CIFAR10 & ResNet-18 & $1.19$ & $0.56$ \\
CIFAR10 & ResNet-50 & $1.12$ & $0.65$ \\
ImageNet100 & CNN & $1.16$ & $0.57$ \\
ImageNet100 & ResNet-18 & $1.12$ & $0.61$ \\
ImageNet100 & ResNet-50 & $1.07$ & $0.67$ \\
\bottomrule
\end{tabular}
\end{table*}

Die dynamische Aktivierungsquantisierung ist in jeder Kombination \emph{langsamer} als FP32 (Faktor $0.21$--$0.67$, entsprechend $1.5$--$4.8\times$ langsamer). Der Grund ist strukturell und in der Diagnose dokumentiert. torchao~0.17s Standard-Filterfunktion für \texttt{Int8DynamicActivationInt8WeightConfig} erfasst ausschließlich \texttt{nn.Linear}-Schichten (\texttt{aten.linear}/\texttt{aten.mm}), nicht \texttt{aten.conv2d}; die einzige Conv2d-fähige Konfiguration, \texttt{IntxWeightOnlyConfig}, dequantisiert bei jedem Vorwärtsdurchlauf zurück auf FP32 und führt eine gewöhnliche FP32-Faltung aus. In dieser torchao-Version existiert schlicht kein optimierter INT8-Faltungskernel. Der moderate Geschwindigkeitsgewinn von \texttt{int8\_weight\_only} ($1.05$--$1.19\times$) stammt entsprechend nicht aus beschleunigter Faltungsberechnung, sondern daraus, dass diese Konfiguration die pro Aufruf anfallenden Aktivierungs-Quantisierungskosten der dynamischen Variante vermeidet, während sie dennoch den Gewichtsspeicher reduziert.
\par\medskip

\textbf{CPU-Pfad (fbgemm).} Der \texttt{fbgemm}-Backend von PyTorchs eager-mode-PTQ implementiert, im Gegensatz zum torchao-GPU-Pfad, echte quantisierte Conv2d- und Linear-Kernel für x86-Hardware mit AVX-512-VNNI (Zielsystem: AMD Ryzen~7 7700X, Zen~4). Die projektinternen Logs bestätigen, dass dieser Pfad für alle sechs Kombinationen ausgeführt wurde. Schichten werden nachweislich in echte quantisierte \texttt{Conv2d}- und \texttt{Linear}-Module konvertiert (\enquote{confirmed quantized}) und auf 50 Kalibrierungs-Batches kalibriert. Ein vollständiger Durchsatz-/Genauigkeits-Sweep mit finaler CSV-Zusammenfassung (analog zu Tabelle~\ref{tab:deployment}) wurde in den verfügbaren Läufen jedoch nicht bis zum Abschluss protokolliert, denn eine \texttt{fbgemm\_summary.csv} existiert in keinem der Ergebnisverzeichnisse. Der CPU-Pfad wird daher qualitativ als mechanistisch tragfähig berichtet (echte INT8-Kernel greifen nachweislich), ohne einen quantitativen Speedup zu behaupten, der sich nicht mit einer Quelle belegen lässt.
\par\medskip

\textbf{Speicherreduktion.} INT8-Gewichte benötigen ein Byte pro Wert gegenüber vier Byte bei FP32, also ein theoretisches Maximum von $4\times$ reiner Gewichtsspeicherreduktion. Ein tatsächlich gemessener Speicherreduktionsfaktor (etwa über nicht-quantisierte Anteile wie Bias- oder BN-Statistiken hinweg) wurde in keinem der verfügbaren Protokolle festgehalten, weshalb nur der theoretische obere Grenzwert und keine gemessene Zahl berichtet wird.
\par\medskip

\textbf{Gesamtbefund.} Der PoT-Bitshift-Vorteil aus \S\ref{sec:fund-quant} übersetzt sich auf der getesteten GPU-Hardware \emph{nicht} in einen gemessenen Laufzeitvorteil. Im Gegenteil, die naheliegendste torchao-Konfiguration ist deutlich langsamer als FP32. Der CPU-Pfad mit echten INT8-Kernen ist mechanistisch der richtige Weg zu einem tatsächlichen Vorteil, konnte im Rahmen dieses Projekts aber nicht bis zu einer abgeschlossenen, quantitativ berichtbaren Messung geführt werden.
\par\medskip

\textbf{Eager-Mode-Referenzmessung.} Ergänzend zum torchao-Pfad liegt für alle sechs Kombinationen eine unabhängige FPS-Messung in PyTorchs eager-mode-PTQ vor, ohne backend-optimierten Kernel (Quelle: \texttt{ptq\_summary.csv}, \texttt{qat\_summary.csv}). Tabelle~\ref{tab:eagerfps} zeigt, dass auch hier PTQ und QAT in jeder Kombination \emph{langsamer} sind als FP32 (Faktor $0.10$--$0.39$), am deutlichsten bei ResNet-50 unter QAT ($0.10$--$0.11\times$). Dies ist ein drittes, methodisch unabhängiges Beispiel für denselben Befund wie \S\ref{sec:phase10} oben: Ohne einen Backend-Kernel, der quantisierte Arithmetik tatsächlich beschleunigt (torchao auf A100, \texttt{fbgemm} auf CPU), bleibt der theoretische PoT-Bitshift-Vorteil ungenutzt, unabhängig davon, mit welcher Software-Schicht gemessen wird.

\begin{table}[t]
\centering
\small
\caption{FP32-, PTQ- und QAT-Durchsatz (Bilder/s) in PyTorchs eager-mode-PTQ, ohne Backend-Kernel. Quelle: \texttt{ptq\_summary.csv}, \texttt{qat\_summary.csv}.}
\label{tab:eagerfps}
\begin{tabular}{llrrrrr}
\toprule
Datensatz & Modell & FP32 & PTQ & QAT & PTQ-Sp. & QAT-Sp. \\
\midrule
CIFAR10 & CNN & 1027.6 & 397.1 & 390.4 & $0.39\times$ & $0.38\times$ \\
CIFAR10 & ResNet-18 & 306.2 & 79.5 & 37.7 & $0.26\times$ & $0.12\times$ \\
CIFAR10 & ResNet-50 & 133.6 & 27.6 & 13.8 & $0.21\times$ & $0.10\times$ \\
ImageNet100 & CNN & 861.6 & 322.0 & 314.4 & $0.37\times$ & $0.36\times$ \\
ImageNet100 & ResNet-18 & 292.9 & 74.9 & 74.7 & $0.26\times$ & $0.25\times$ \\
ImageNet100 & ResNet-50 & 127.8 & 26.5 & 13.7 & $0.21\times$ & $0.11\times$ \\
\bottomrule
\end{tabular}
\end{table}

\subsection{Accuracy-vs-Loss-Robustheits-Check}
\label{sec:robustness}

Um zu prüfen, ob die in \S\ref{sec:phase4} berichteten Befunde artefaktisch von der Wahl der Loss-basierten Schaden-Definition (\S\ref{sec:isolation}) abhängen, wird die identische Analyse mit genauigkeitsbasiertem Schaden wiederholt (\texttt{weight\_ablation\_canonical\_v2.csv} und dessen ImageNet100-Variante, \texttt{weight\_ablation\_canonical\_correlation\_v2.csv}). Die qualitativen Kernaussagen bleiben unverändert. Die wahre Top-Damage-Schicht ist in allen 6 von 6 PTQ-Kombinationen identisch zwischen der Loss- und der Genauigkeits-basierten Definition (\texttt{conv1}, \texttt{conv1}, \texttt{layer1.1.conv1}, \texttt{conv1}, \texttt{fc}, \texttt{conv1} für die sechs Kombinationen aus Tabelle~\ref{tab:rank}). Kein einziges der 18 Score$\times$Kombination-Korrelationsvorzeichen kehrt sich zwischen den beiden Definitionen um.
\par\medskip

Was sich verschiebt, ist die Signifikanz-Zusammensetzung. Unter genauigkeitsbasiertem Schaden sind 3 von 18 Korrelationen bei $p<0.05$ signifikant (\texttt{cnn}/CIFAR10 $S_{\mathrm{raw}}$, ResNet-50/ImageNet100 $S_{\mathrm{raw}}$, \texttt{cnn}/ImageNet100 $S_{\mathrm{pert}}$), unter loss-basiertem Schaden 5 von 18 (\S\ref{sec:phase4}), teils dieselben, teils andere Score$\times$Kombination-Paare. Dass mehr Korrelationen unter dem Loss-Ziel signifikant werden, ist konsistent mit der in \S\ref{sec:isolation} genannten Motivation, den Loss als die für Gleichung~\eqref{eq:taylor} theoretisch passendere, kontinuierlichere Zielgröße zu verwenden.
\par\medskip

Abbildung~\ref{fig:robustness} macht diese Übereinstimmung sichtbar, indem sie für alle 18 Score$\times$Kombination-Paare die loss-basierte gegen die genauigkeitsbasierte Spearman-$\rho$ aufträgt. Alle Punkte liegen nahe der Diagonalen $y=x$ und, entscheidend, auf derselben Seite der Null-Linie wie ihr loss-basiertes Gegenstück, was die in \S\ref{sec:phase4} berichtete qualitative Konsistenz noch einmal direkt zeigt statt sie nur zusammenzufassen. Insgesamt stützt dieser Robustheits-Check die Kernaussage aus \S\ref{sec:phase4} unabhängig von der gewählten Schaden-Metrik, eine qualitative Bestätigung, dass Loss die richtige Zielgröße für diese Analyse ist (\S\ref{sec:isolation}), ohne dass Genauigkeit dafür jemals als Alternative geführt wurde.

\begin{figure}[t]
\centering
\includegraphics[width=0.85\linewidth]{figures/fig_05_acc_loss_comparison.pdf}
\caption{Spearman-$\rho$ je der 18 Score$\times$Kombination-Paare, loss-basiert (x-Achse, \S\ref{sec:phase4}) gegen genauigkeitsbasiert (y-Achse, dieser Abschnitt). Graue Diagonale $=y=x$. Erzeugt von \texttt{notebooks/13\_accuracy\_loss\_robustness.ipynb}.}
\label{fig:robustness}
\end{figure}